%%=============================================================================
%% Samenvatting
%%=============================================================================

% TODO: De "abstract" of samenvatting is een kernachtige (~ 1 blz. voor een
% thesis) synthese van het document.
%
% Een goede abstract biedt een kernachtig antwoord op volgende vragen:
%
% 1. Waarover gaat de bachelorproef?
% 2. Waarom heb je er over geschreven?
% 3. Hoe heb je het onderzoek uitgevoerd?
% 4. Wat waren de resultaten? Wat blijkt uit je onderzoek?
% 5. Wat betekenen je resultaten? Wat is de relevantie voor het werkveld?
%
% Daarom bestaat een abstract uit volgende componenten:
%
% - inleiding + kaderen thema
% - probleemstelling
% - (centrale) onderzoeksvraag
% - onderzoeksdoelstelling
% - methodologie
% - resultaten (beperk tot de belangrijkste, relevant voor de onderzoeksvraag)
% - conclusies, aanbevelingen, beperkingen
%
% LET OP! Een samenvatting is GEEN voorwoord!

%%---------- Nederlandse samenvatting -----------------------------------------
%
% TODO: Als je je bachelorproef in het Engels schrijft, moet je eerst een
% Nederlandse samenvatting invoegen. Haal daarvoor onderstaande code uit
% commentaar.
% Wie zijn bachelorproef in het Nederlands schrijft, kan dit negeren, de inhoud
% wordt niet in het document ingevoegd.

\IfLanguageName{english}{%
\selectlanguage{dutch}
\chapter*{Samenvatting}
\selectlanguage{english}
}{}

%%---------- Samenvatting -----------------------------------------------------
% De samenvatting in de hoofdtaal van het document

\chapter*{\IfLanguageName{dutch}{Samenvatting}{Abstract}}

De moderne voetbalwereld beschikt over enorme hoeveelheden aan data en er wordt altijd gezocht naar nieuwe manieren om deze data te gebruiken om de prestaties van de club te verbeteren. Zo zoekt ook KAA Gent naar allerhande manieren om deze data te gebruiken, een domein waar er nog ruimte is voor verbetering is in het scoutingsdepartement, meer specifiek in het onderzoek naar het belang van fysieke prestaties van spelers bij het behalen van overwinningen.
\newline \newline
De huidige methode van het scoren van de fysieke prestaties is gebaseerd op het arbitrair inschatten van het belang van bepaalde fysieke statistieken. Dit is echter een subjectieve en geen objectieve inschatting van de data. Een manier om deze zaken objectiever te beoordelen is het te benaderen vanuit een machine learning-perspectief.
\newline \newline
De centrale onderzoeksvraag van dit onderzoek luidt: \emph{Hoe kan een objectieve en reproduceerbare fysieke spelersscore ontwikkeld worden op basis van wedstrijddata met behulp van machine learning-technieken?} Hierbij wordt onderzocht welke methode het beste is om het belang van bepaalde statistieken te bepalen, alsook hoe deze benadering past in de huidige scoutingstructuur.
\newline \newline
Het hoofddoel van het onderzoek is het construeren van een fysieke spelersscore gebaseerd op de inzichten van een machine learning-model dat wedstrijdresultaten voorspelt.
\newline \newline
De methodologie die in dit onderzoek gehanteerd wordt begint bij een uitgebreide literatuurstudie, waarna een proof-of-concept gecreëerd wordt en wordt toegevoegd aan een alternatieve versie van de huidige scoutingstool binnen de club.
\newline \newline
Het uiteindelijke resultaat van dit onderzoek toont aan dat het mogelijk is om een fysieke spelersscore gebruikmakend van machine learning-technieken te berekenen. De gecreëerde modellen hebben echter slechts beperkte voorspellingskracht. De oorzaak hiervoor ligt hoogstwaarschijnlijk aan het feit dat een voetbaluitslag niet louter beslist wordt door de fysieke prestaties van spelers op het veld maar ook aan veel andere factoren, zoals technische statistieken, tactische beslissingen en zelfs weersomstandigheden.
\newline \newline
De conclusie van dit onderzoek is dat fysieke prestaties een beperkte invloed hebben op een wedstrijd. Het is echter wel mogelijk om een fysieke score te berekenen met machine learning-technieken en met verder onderzoek is het wellicht mogelijk om een nog betere en objectievere fysieke spelersscore te ontwikkelen.